% !TEX root = ../notes.tex

\documentclass[../notes.tex]{subfiles}

\begin{document}

\section{April 27}
Here we go.

\subsection{Condensed Fixed Points}
Recall that we let $\Gamma_n$ denote the Honda formal group law of height $n$, and then for a perfect field $k$ of positive characteristic $p$, the Morava $E$-theory $E_n\coloneqq E(k,\Gamma_n)$. Then $E_n$ had an action by $\mathbb G_n\coloneqq\OO_n^\times\rtimes\op{Gal}(k/\FF_p)$, where $\OO_n\coloneqq W(k)\langle S\rangle/\left(S^n-p,Sa-a^\sigma S\right)$.

It will turn out that $L_{K(n)}\mathbb S=E_n^{h\mathbb G_n}$, and we will explain what this means using the language of condensed mathematics, which provides a good framework for giving algebraic objects the action of a topological group. To understand why this framework is useful, we note that it will also turn out that $L_{K(n-1)}L_{K(n)}\mathbb S=(L_{K(n-1)}E_n)^{h\mathbb G_n}$. However, it is not so clear how we would describe the ``continuous action'' by $\mathbb G_n$ on $L_{K(a)}E_n$ for any $a$. For example, the homotopy ring is
\[\pi_*L_{K(a)}E_n=\left(W(k)[[v_1,\ldots,v_{n-1}]][u,1/u][1/v_a]\right)^\land_{(p,v_1,\ldots,v_{a-1})},\]
but it is not so obvious how to give this thing a topology which will successfully make the multiplication continuous. Instead, it will turn out that we can descend the power series ring $W(k)[[v_1,\ldots,v_{n-1}]][u,1/u]$ from a topological ring to a condensed ring, and then the remaining localizations and completion preserve condensed rings.

With this in mind, we now want to view $E_n$ as a condensed spectrum.
\begin{remark}
	There are two ways to make $E_n$ into a condensed spectrum.
	\begin{itemize}
		\item We could write
		\[E_n=\lim\left(E_n/\left(p^{a_0},v_1^{a_1},\ldots,v_{n-1}^{a_{n-1}}\right)\right).\]
		\item We could use \Cref{thm:lurie} to build a condensed spectrum by the functor $S\mapsto E(C(S,k),\Gamma_n)$. Note that this outputs to $\mathbb E_\infty$-rings. Here, $C(S,k)$ consists of the continuous functions from $S$ to $k$.
	\end{itemize}
	It turns out that these two constructions give isomorphic condensed spectra. Notably, $\mathbb G_n$ as a condensed group acts functorially on the second construction, so we receive the action of a condensed group 
\end{remark}
\begin{remark}
	More generally, we can view $E$ as a functor from light condensed rings $R$ to light condensed $\mathbb E_\infty$-spaces. In particular, it should take as input a light condensed ring $R$ and a $p$-divisible group over $R(\mathrm{pt})$. In order to get a well-behaved (even) answer, one should require $(R(\mathrm{pt}),\Gamma)$ is Luriean and that the internal maps $R(\mathrm{pt})\to R(S)$ are flat and relatively perfect for each $S$.
\end{remark}
And now we can take fixed points.
\begin{example}
	We can now make sense of $E_n^{h\mathbb G_n}$. By unwinding the definitions, we are interested in computing the limit of the hypercover
	\[\lim E(C(\mathbb G_n^\bullet,k),\Gamma_n).\]
	Roughly speaking, we are decomposing this limit into a sifted limit of products.
\end{example}
\begin{remark}
	Let's explain the formula $\lim E(C(\mathbb G_n^\bullet,k),\Gamma_n)$. For example, if a group $G$ acts on a spectrum $X$, then $X^{hG}=\lim_{\mathrm BG}X$. Then we note that $\mathrm BG$ can be found as the colimit of the simplicial object $\cdots\to G\times G\to G\to*$. Taking the limit back discovers that
	\[X^{hG}=\lim G^\bullet\times X.\]
\end{remark}
\begin{example}
	Note that $\pi_0E_h=W(k)[[u_1,\ldots,u_{n-1}]]$ admits a deformation of $\widetilde\Gamma_n$. It turns out that
	\[L_{K(a)}E_h=E\left(\pi_0E_h/(p,u_1,\ldots,u_{a-1})[1/u_a],\widetilde\Gamma_h\right).\]
\end{example}

\subsection{Application to Chromatic Splitting}
Now, for the chromatic splitting conjecture, we are interested in $L_{K(h-1)}L_{K(h)}\mathbb S$. We might want to be able to see this object on $E$-theory.
\begin{notation}
	Let $k$ be a field containing both $\FF_{p^{h-1}}$ and $\FF_{p^h}$. Then we define
	\[\mathbb B=E_{h-1}\otimes_{W(k)}L_{K(h-1)}E_h.\]
\end{notation}
\begin{remark}
	By construction, $\mathbb B$ admits a continuous action by $\mathbb G_{h-1,n}\coloneqq\mathbb G_{n-1}\times_{\op{Gal}(k/\FF_p)}\mathbb G_n$. We remark that taking fixed points by one of these $\mathbb G_n$ actually retracts to a finite limit (this means that $\mathbb G_n$ has finite cohomological dimension). It then turns out that $\mathbb B^{h\mathbb G_{h-1,n}}=L_{K(h-1)}L_{K(n)}\mathbb S$, which we can see basically by taking the fixed points on each factor separately (which is possible because the limits are finite).
\end{remark}
All interesting even rings should be found via $E$-theory. Here is how we do this for $\mathbb B$.
\begin{notation}
	Set $K\coloneqq\pi_0\left(L_{K(h-1)}E_h/(p,v_1,\ldots,v_{h-2})\right)$, which can be described as $\FF_p((u_{h-1}))$. It turns out that there is a Galois extension $L$ of $K$ which is the smallest field extension of $K$ over which there is an isomorphism $\Gamma_{h-1}\cong(\widetilde\Gamma_h)^\circ$.
\end{notation}
\begin{remark}
	It turns out $(\widetilde\Gamma_h)^\circ$ has height $h-1$, which one can more or less see on the level of the formal groups; notably, taking the connected component can lower the height! By general principles, one can show that the scheme $\op{Isom}_k\left(\Gamma_{h-1},(\widetilde\Gamma_h)^\circ\right)$ is a torsor over $K$ and thus \'etale over $K$. It will turn out that this scheme is actually connected, so it is a field extension of $K$, which will be $\Spec L$.
\end{remark}
\begin{remark}
	It is a theorem of Torii that $\mathbb B$ is $E(L,\Gamma_{h-1})$.
\end{remark}
\begin{remark}
	It turns out that the $u_{h-1}$-adic completion $\widehat L$ is isomorphic to $\FF_p((t^{1/p^\infty}))$. Jeremy Hahn has been told that this is a useless isomorphism.
\end{remark}
\begin{notation}
	We define $\widehat{\mathbb B}\coloneqq E(\widehat L,\Gamma_{h-1})$.
\end{notation}
Now, there is a composite
\[E_{h-1}\to\mathbb B\to\widehat{\mathbb B}.\]
Indeed, this is just the extension composite $E(k,\Gamma_{h-1})\to E(L,\Gamma_{h-1})\to E(\widehat L,\Gamma_{h-1})$. One can take some invariants to produce a sequence
\[E_{h-1}^{h\ZZ_p^\times}\to\mathbb B^{h\mathbb G_h}\to\widehat{\mathbb B}^{h\mathbb G_h}\]
which is $\OO_{h-1}^{\times,1}$-equivariant, where $\OO_{h-1}^{\times,1}$ is the kernel of the projection $\OO_{h-1}^\times\to\ZZ_p^\times$.

Here is a big theorem which is currently work in progress.
\begin{theorem} \label{thm:easy-chromatic-splitting}
	Take $p$ large. Then the composite $E_{h-1}^{h\mathbb Z_p^\times}\to\widehat{\mathbb B}^{h\mathbb G_h}$ is an isomorphism.
\end{theorem}
\begin{remark}
	Taking $\OO_{h-1}^{\times,1}$-fixed points, we exhaust our group actions, and we receive a section of the map $L_{K(h-1)}\mathbb S\to L_{K(h-1)K(h)}\mathbb S$, which is a basic consequence of \Cref{conj:chromatic-splitting}.
\end{remark}
This result arises from some calculation.
\begin{theorem}
	Let $\mathbb G_h^1$ be the kernel of some suitable determinant $\det\colon\mathbb G_h\to\ZZ_p^\times$. Then
	\[\mathrm H^*_{\mathrm{cts}}\left(\mathbb G_h^1;\widehat L\right)\cong k\oplus\Sigma^{-h}k\]
	with residual nontrivial $\ZZ_p^\times$-action on $\Sigma^{-h}k$.
\end{theorem}
\begin{remark}
	Let's explain how this implies \Cref{thm:easy-chromatic-splitting}. To start, it is enough to check that our map $E_{h-1}^{h\mathbb Z_p^\times}\to\widehat{\mathbb B}^{h\mathbb G_h}$ is an isomorphism on homotopy groups, and it is even enough to check this after modding out by $(p,v_1,\ldots,v_{n-2})$ because everything is local (which means that we can check after taking $K(n)$-homology). Well,
	\[\pi_*E_{h-1}=W(k)[[u_1,\ldots,u_{h-2}]][u,1/u],\]
	and taking the quotient leaves us with $k$, so taking the fixed points yields $k\oplus\Sigma^{-1}k$. On the other hand, the homotopy groups of $\pi_*(\widehat{\mathbb B}^{h\mathbb G_h^1})$ are computed with some spectral sequence (for large primes) via this continuous cohomology. Taking fixed points by $\ZZ_p^\times$ then recovers $k\oplus\Sigma^{-1}k$. Note that the map in homotopy groups is automatically an isomorphism because it is basically induced by the identity on $k$ (which is the only ring map $k\to k$ anyway) before taking fixed points by $\ZZ_p^\times$.
\end{remark}

\end{document}